\documentclass[preview]{standalone}

\usepackage{amsmath}
\usepackage{amssymb}
\usepackage{stellar}
\usepackage{definitions}

\begin{document}

\id{psicologia-se-valutazione-personalita}
\genpage

\section{Sé e valutazione della personalità}

\begin{snippet}{james-se}
    William James descrisse il sé come una struttura mentale dinamica capace di
    motivare, interpretare, organizzare, mediare e regolare comportamenti e
    processi intrapersonali e interpersonali.
\end{snippet}

\begin{snippet}{componenti-se-james}
    James individuò tre componenti dell'esperienza di sé:
    \begin{itemize}
        \item \emph{me materiale}, cioè il corpo e gli oggetti fisici che lo
        circondano;
        \item \emph{me spirituale}, cioè il sé che monitora pensieri e
        sentimenti intimi;
        \item \emph{me sociale}, cioè la consapevolezza di come gli altri ci
        vedono.
    \end{itemize}
\end{snippet}

\begin{snippetdefinition}{schema-se-definition}{Schema di sé}
    Uno \emph{schema di sé} è una struttura cognitiva che organizza le
    informazioni relative a sé stessi e guida il modo in cui tali informazioni
    vengono elaborate.
\end{snippetdefinition}

\begin{snippetdefinition}{autostima-definition}{Autostima}
    L'\emph{autostima} è una valutazione generale del sé, influenzata
    dall'interazione tra fattori genetici ed esperienze ambientali.
\end{snippetdefinition}

\begin{snippetdefinition}{autoaccrescimento-definition}{Autoaccrescimento}
    L'\emph{autoaccrescimento} comprende strategie finalizzate a interpretare le
    proprie azioni e i propri comportamenti in modo costantemente positivo.
\end{snippetdefinition}

\begin{snippetdefinition}{strategie-autolesive-definition}{Strategie autolesive}
    Le \emph{strategie autolesive} sono processi che, in previsione di un
    possibile fallimento, portano a produrre comportamenti o spiegazioni capaci
    di ridurre la probabilità di attribuire l'insuccesso alle proprie
    insufficienti capacità.
\end{snippetdefinition}

\begin{snippetdefinition}{costrutto-indipendente-se-definition}{Costrutto indipendente del sé}
    Il \emph{costrutto indipendente del sé}, tipico delle culture
    individualistiche, concepisce il sé come individuo il cui comportamento è
    organizzato principalmente attorno ai propri pensieri, sentimenti e azioni.
\end{snippetdefinition}

\begin{snippetdefinition}{costrutto-interdipendente-se-definition}{Costrutto interdipendente del sé}
    Il \emph{costrutto interdipendente del sé}, tipico delle culture centrate
    sulla collettività, concepisce il sé come parte di una relazione sociale
    inclusiva, in cui il comportamento è organizzato anche dai pensieri,
    sentimenti e azioni altrui.
\end{snippetdefinition}

\begin{snippet}{teorie-personalita-confronto}
    Le principali teorie della personalità possono essere confrontate lungo
    alcune opposizioni: ereditarietà e ambiente, apprendimento e leggi innate,
    enfasi sul passato o sul futuro, processi consci e inconsci, disposizioni
    interne e situazioni esterne.
\end{snippet}

\begin{snippet}{valutazione-personalita-introduzione}
    La valutazione della personalità parte da due assunti: esistono
    caratteristiche individuali che rendono coerente il comportamento di una
    persona, e tali caratteristiche possono essere valutate o misurate. I test
    devono rispettare criteri di affidabilità e validità.
\end{snippet}

\begin{snippetdefinition}{inventario-personalita-definition}{Inventario della personalità}
    Un \emph{inventario della personalità} è un test oggettivo basato su un
    questionario self-report, composto da item relativi a pensieri, sentimenti e
    azioni.
\end{snippetdefinition}

\begin{snippetexample}{mmpi-neo-pi-example}{MMPI e NEO-PI}
    Il \emph{Minnesota Multiphasic Personality Inventory}, o \emph{MMPI}, è
    usato anche per diagnosi psichiatriche e include scale cliniche. Il
    \emph{NEO-PI} è invece progettato per valutare caratteristiche di
    popolazioni adulte non cliniche ed è fondato sul modello a cinque fattori.
\end{snippetexample}

\begin{snippetdefinition}{test-proiettivo-definition}{Test proiettivo}
    Un \emph{test proiettivo} è un test che presenta stimoli ambigui, come
    macchie, immagini incomplete o disegni interpretabili in più modi, chiedendo
    alla persona di descriverli, completarli o costruire storie a partire da
    essi.
\end{snippetdefinition}

\begin{snippetdefinition}{test-rorschach-definition}{Test di Rorschach}
    Il \emph{test di Rorschach} è un test proiettivo in cui vengono mostrate
    macchie d'inchiostro, alcune in bianco e nero e altre a colori, chiedendo
    alla persona che cosa potrebbero rappresentare.
\end{snippetdefinition}

\begin{snippetdefinition}{test-appercezione-tematica-definition}{Test di appercezione tematica}
    Il \emph{test di appercezione tematica}, o \emph{TAT}, è un test proiettivo
    in cui vengono mostrate immagini ambigue e viene chiesto al soggetto di
    produrre storie sui personaggi, sugli eventi e sul loro possibile esito.
\end{snippetdefinition}

\begin{snippet}{uso-test-personalita}
    Ogni tipo di test fornisce informazioni specifiche sulla personalità. Per
    questo motivo, nelle valutazioni cliniche i professionisti usano spesso una
    combinazione di test oggettivi e proiettivi.
\end{snippet}

\end{document}
